\begin{proof}[Proof of Corollary~\ref{cor:no-glob}]
(1) A finite language with two words is neither a singleton nor unbounded, so
Lemma~\ref{lem:dichotomy} excludes it. (2) $\varepsilon\notin\Sigma^{+}$ forces
$m(g)\ge 1$, since $m(g)=0$ makes $g$ a string of stars and $\varepsilon\in L(g)$; and
$\Sigma^{+}$ contains words of length $1$, so $m(g)\le 1$. Hence $m(g)=1$: $g$ has a
single literal character $c$ and at least one star. By Lemma~\ref{lem:dichotomy} every
word of $L(g)$ contains $c$, while $\Sigma^{+}$ contains the one-character words of
$\Sigma\setminus\{c\}$, which is non-empty. (3) The same argument gives $m(g)=1$ with one
literal $c$, which must lie in $A$ since every word of $L(g)$ contains it. A single
literal and at least one star leaves three shapes: $c{*}$, ${*}c$ and ${*}c{*}$. In the
first and third $L(g)$ contains $c\sigma$ for every $\sigma\in\Sigma$; in the second it
contains $\sigma c$ for every $\sigma$. Either way it contains a word using a
$\sigma\notin A$, which $A^{+}$ does not, and $\Sigma\setminus A\neq\emptyset$.

Taking $c\sigma$ in every case, as an earlier version of this proof did, is wrong for
$g={*}c$: that language has no word beginning with $c$ unless $c$ repeats.
\end{proof}
